\section{Strong irreducibility and the marked rank-three automorphy implication}
\label{sec:markedaut-v64}

% BEGIN VMA1
\paragraph{VMA1. The absolute ordinary point on the prescribed host.}
\label{claim:VMA1}
Fix any $K$ with the source Gaussian CM condition $\mathrm{IsQi}$,
an algebraic closure $\Omega/\mathbb Q$, a prime $\ell$, the
comparison $\iota:\overline{\mathbb Q}_\ell\simeq\mathbb C$,
$c:\mathrm{Coef}(\ell)$, $N\ge1$, $e,\psi,o_1$ with
$\mathrm{HostData}$, and $S,(A_v),F$ with $\mathrm{IsHostFrame}$.
Put $O=c.O$, $k=O/\mathfrak m_O$, $E=\overline{\mathbb Q}_\ell$,
$\Delta=G_{\mathbb Q}$ on this $\Omega$, $\Gamma=G_S$ and
$p:\Delta\twoheadrightarrow\Gamma$. The pullbacks of the frame
characters are $\eta_4$ and $\Psi_0=\chi_\ell^e\psi_{\rm Gal}$.
Write $H=\ker\eta_4$, choose the frame's complex conjugation lift
$\sigma$ with $p\sigma=F.c$, and put $B=A_{v_\ell}$. Then
$\sigma^2=1$, $\eta_4(\sigma)=-1$, $\Psi_0(\sigma)=1$ by
complex conjugation and the given parity, and $v_B(\ell)<1$ in
multiplicative valuation notation. AFE10 supplies
$x\in D_B\cap H$ with $\Psi_0(x)^2-1\in O^\times$, including
at $\ell=3$. Thus a marked element exists already. In the
conditional joins below $g_*$ is the particular absolute marking
supplied jointly with the arithmetic ends, and $g_1=p(g_*)$.
In each case the identities $\eta_4(g_*)=1$ and
$\Psi_0(g_*)^2-1\in O^\times$ give $F.\mathrm{Marked}(g_1)$.

For every integral local coefficient algebra $A$ on which the point
conditions are defined, a marked point $r':\Gamma\to\mathrm{GG}_3(A)$
at $g_1$ gives the original absolute $\mathrm{OrdPt}$ at
$(\sigma,g_*,B)$ by pullback. The sign and residual congruences
pull back entrywise. The reductions are locally constant since
$p$ is continuous. The normalization at $\sigma,g_*$ follows from
$p\sigma=F.c$ and $p(\sigma g_*^{-1}\sigma)=F.c g_1^{-1}F.c$.
The prescribed upper triangular flag at $B$, including its inertial
diagonal $(\Psi_0,1,\Psi_0^{-1})$ on $H$, is the point's own
$\mathrm{CondT}$ after identifying the actual $D_B,I_B$.

For the away condition take any inertia element at a valuation
away from $\ell$ on which $\eta_4=\Psi_0=1$. The trivial
valuation has trivial inertia. Outside $S$ the quotient kills
the element. Inside $S$, conjugate the valuation to the chosen
$A_v$. At $v\in S_b$ the local diagonal condition makes all
diagonal entries one, hence the whole element one in its diagonal
basis. At the other $v\ne v_\ell$ the element is in
$I_v\cap\ker\eta_v\cap\ker\bar\Psi_v$, on which the point is
trivial. Conjugating back preserves the identity in $\mathrm{GG}_3$.
This verifies all five fields of the absolute ordinary point for the
same datum, in both residual and characteristic-zero specializations.
% END VMA1

% BEGIN VMA2
\paragraph{VMA2. Finite weight zero from the absolute finite end.}
\label{claim:VMA2}
Fix $g_*,g_1$ as in VMA1. Assume its $\mathrm{FinEnd}$: there is
one finite-index $V\le\Delta$ such that, over every complete DVR
$D$ with finite residue field and local $O$-structure annihilating
$\mathfrak m_O$, each absolute ordinary point on this datum has
the base characteristic polynomial on $V\cap H$. Let $R,u$ be
\emph{any} universal marked pair at $g_1$.
The host frame has $F.\mathrm{Ok}$ by the frame verifications in
CVI2: the decomposition separator is AFE10 for every $v\in T_b$,
away from $T_b$ the ramified part has order prime to $\ell$, and
the continuity, parity and Frobenius conditions follow from the
defining identifications of this host. These verifications use
$\mathrm{HostData}$ and $\mathrm{IsHostFrame}$ alone. The source
groups are profinite (the quotient kernel is closed), and $S$ is
finite. NMD3's marked-entry argument therefore gives a finite
adapted weighted presentation of this $R$ with variables
$z_i=u(h_i)_{a_i b_i}-(r_0(h_i))_{a_i b_i}$, $h_i\in\ker\eta$,
of weights $w_i=b_i-a_i$. They lie in $\mathfrak m_R$, and its
weight-zero component is $\mathrm{wpart}(O,u,0)$.

Take any local $O$-map $R\to D$ of the kind tested in UCE9.
It sends the universal point to an integral marked point:
congruences and local bases pass along the map, and
$\mathfrak m_R^n$ maps into $\mathfrak m_D^n$, giving continuity.
VMA1 gives its absolute ordinary realization. Since
$O\to D$ factors through $k$, the base character in the
$\mathrm{FinEnd}$ equality here is precisely the reduction
$\bar\Psi_0$. UCE10--11 apply with the original finite $k$,
$2\in k^\times$, the marked element $g_*$, the original normalization
and the given $V$. They yield one triangular direction on $H$ and
equality of all diagonal entries with the base. Explicitly, a positive
power of every element of $H$ lies in $V$; its eigenvalues are thus
algebraic over $k$. Adjoining an algebraic element of the local domain
gives a field over $k$, so a characteristic-polynomial coefficient
congruent to its base value in $k$ equals that base value. The marked
diagonal projectors and UCE11's trace-pairing argument then exclude
both a three-dimensional irreducible block and an irreducible outer
two-dimensional block. This works in characteristic three as well.

The map $p$ sends $H$ onto $\ker\eta$. The test map kills the $z_i$
of weight zero, and kills every $z_i$ on one of the two strict sides
of zero: off the diagonal the base entries vanish. Thus each
nonconstant monomial of weight zero vanishes at every such test
map. By the DVR-curve criterion UCE9 all these test maps imply
$\mathrm{wpart}(O,u,0)$ is a finite $O$-module. In particular
this conclusion is integral and allows nilpotents in $R$.
% END VMA2

% BEGIN VMA3
\paragraph{VMA3. The marked eigenvalue test on every open subgroup.}
\label{claim:VMA3}
Keep any mark $g_1$. Let $r:\Gamma\to\mathrm{GG}_3(E)$ be an
original characteristic-zero marked point of Type I. It has an
integral marked representative $r'$ over a finite DVR $A\subset E$
over $O$, related by conjugation by $T(t)=\operatorname{diag}(t,1,t^{-1})$
for one $t\in E^\times$. On $H$ the matrix of $r\circ p$ is an
ordinary representation $\rho$, as the sign there is positive.
It is irreducible because $p(H)=\ker\eta$ and Type I means
irreducibility on that kernel. The matrices of $r\circ p$ are
continuous: the local-constant reductions of $r'$ give continuity
in its finite coefficient field, and each matrix of $r$ is obtained
by fixed torus conjugation. For every open subgroup $U$ of
$\Delta$ contained in $H$, the restriction $\rho|_U$ is irreducible.

To prove this, take its finite-index normal core $U_0$ in $\Delta$.
For the chosen $B$ over $\ell$, take $s\in I_B$ with
$\chi_\ell(s)=1+\ell$. Indeed the cyclotomic extensions
$\mathbb Q_\ell(\mu_{\ell^n})$ are totally ramified, by the
Eisenstein polynomial for $\zeta_{\ell^n}-1$; the images on
inertia at each finite level are full, and compactness lifts the
compatible units to a single element. With
$m=[\Delta:U_0]\,|(\mathbb Z/N)^\times|$ we have
$u=s^m\in U_0\cap I_B$ and $\Psi_0(u)=(1+\ell)^{em}$.
The integer $em$ is positive. The three eigenvalues
$\Psi_0(u),1,\Psi_0(u)^{-1}$ in the ordinary basis are distinct.

Suppose $\rho|_{U_0}$ is reducible. A simple $U_0$-subspace exists
by minimal dimension. Its $H$-conjugates are simple and span the
whole irreducible $H$-space, so the restriction is semisimple.
The isotypic components form one $H$-orbit (an orbit sum is
$H$-stable). Their common dimension divides three. A reducible
restriction is therefore scalar, or consists of three distinct lines
permuted transitively: in the single isotypic case the multiplicity
times the dimension of the simple module is three. The scalar case
contradicts the three eigenvalues. The three lines in the second case
are the eigenlines of $u$. Its triangular matrix with distinct
diagonal entries can be diagonalized by an upper triangular basis
change, by choosing the eigenvector for the $i$th diagonal entry
in the $i$th flag step. Hence $H$ acts by monomial matrices in
this basis and the separator $x\in D_B\cap H$ of VMA1 remains
upper triangular there. Its permutation is the identity. Here the
separator is needed only in the decomposition group, and the
regular element $u$ is in inertia.

A transitive subgroup of $S_3$ contains a three-cycle, since its
order is divisible by three. Choose $g\in H$ mapping to this
cycle. Then $\operatorname{tr}\rho(gx^n)=0$ for $n=0,1,2$.
The equality also holds for the integral representative, by
conjugation on $H$. In its residue field put
$a=\overline{\Psi_0(g)}$, $b=\overline{\Psi_0(x)}$. Residual
diagonality yields $ab^n+1+(ab^n)^{-1}=0$ at these three indices.
Thus $a,ab,ab^2$ are roots of $X^2+X+1$. They are three distinct
elements since $b^2-1$ is a unit before reduction. This contradiction
proves irreducibility on $U_0$, hence on $U$.
% END VMA3

% BEGIN VMA4
\paragraph{VMA4. Irreducibility on the full cyclotomic kernel.}
\label{claim:VMA4}
For the same point, put $J=H\cap\ker\chi_\ell$. Then
$\rho|_J$ is irreducible over $E$. This subgroup can have infinite
index. Suppose the restriction is reducible. The sum of all
$H$-conjugates of a simple $J$-subspace is $H$-stable and is
all of $E^3$ by irreducibility. A finite direct sum of simple
subspaces can be extracted from this sum by successively adjoining
a simple subspace that increases the span. Thus $\rho|_J$ is
semisimple. It admits a nonzero proper equivariant projection $P$.
In $\operatorname{End}_J(E^3)$ the element
$P-(\operatorname{tr}P/3)I$ is traceless and nonzero.

The traceless part of this commutant is stable under conjugation
by $H$. Such conjugations commute on it, since $H/J$ is abelian.
A commuting family on a nonzero finite-dimensional space over the
algebraically closed field $E$ has a common eigenvector: if any
member is nonscalar, restrict to a nonzero proper eigenspace and
induct on dimension. Hence a nonzero traceless matrix $M$ satisfies
\[
       \rho(h)M\rho(h)^{-1}=\varphi(h)M\qquad (h\in H)
\]
for a character $\varphi:H\to E^\times$. The kernel of $M$ is
$H$-stable, so $M$ is invertible. Taking determinants gives
$\varphi^3=1$. Its kernel is the commuting locus of $M$ and the
continuous $\rho(h)$, hence is closed. A closed finite-index
subgroup is open (its complement is a finite union of closed
cosets), and $H$ itself is open in $\Delta$. By VMA3 the kernel
acts irreducibly. The eigenspace of any eigenvalue of $M$ is stable
under that kernel and nonzero, hence all of $E^3$. This makes
$M$ scalar, contrary to its nonzero trace-zero condition in
characteristic zero. The argument also applies when $\ell=3$.
% END VMA4

% BEGIN VMA5
\paragraph{VMA5. The reduced and full marked assertions are equivalent.}
\label{claim:VMA5}
The reduced source assertion $\mathrm{ThmS6Core}_s$ supplies,
after the host, frame and $\iota$ are fixed, one marked $g_1$.
For every universal $R,u$ at this mark it gives finite weight zero;
for every $O$-algebra map $\xi:S(u)\to E$ from the graded
entry subalgebra and every group map $r:\Gamma\to\mathrm{GG}_3(E)$
satisfying $\mathrm{IsPtHom}(u,\xi,r)$ and Type I it gives
$\mathrm{AutOnAlt}(K,\iota,-1,(r\circ p)^{\rm mat})$.
With exactly these quantifiers it implies $\mathrm{ThmS6}_s$.

Indeed keep that mark and the supplied finiteness for the given pair.
The characteristic-zero realization VSCM8 applies to this universal
ring with its adapted presentation from NMD3. Its hypotheses are
verified on this instance by the profinite quotient, $F.\mathrm{Ok}$,
the mark and the supplied weight-zero finiteness. Concretely the
weight-zero coordinates land integrally, the nonconstant mixed-sign
monomials of weight zero land in the maximal ideal, and one scalar
$t\in E^\times$ places all the rescaled marked coordinate differences in the maximal ideal
of a single finite coefficient DVR $A$. Evaluation of the homogeneous
relations extends to the completed local ring. The resulting
integral $r_A$ satisfies
\[
     \mathrm{push}_{A\hookrightarrow E} r_A
                    =\mathrm{conj}_{T(t)} r.
\]
The entry and sign identities in $\mathrm{IsPtHom}$ give this equality
for the particular $r$ in the quantifier. This is
$F.\mathrm{IsCharZeroPt}(g_1,r)$. VMA3--4 now provide the two
additional fields $\mathrm{opens}$ and $\mathrm{cyc}$ of $B'_s$ on
the actual absolute Galois subgroups. Its third field is precisely
the supplied $\mathrm{AutOnAlt}$ for this same $\xi,r$ and comparison
$\iota$. This finishes the full implication. Conversely, retaining the
given mark and finite module and projecting $B'_s$ to its third field
gives $\mathrm{ThmS6Core}_s$ immediately.
% END VMA5

% BEGIN VMA6
\paragraph{VMA6. An arithmetic end with one mark suffices.}
\label{claim:VMA6}
The original absolute $\mathrm{HostEndAlt}$ implies
$\mathrm{ThmS6}_s$. To see this, choose the specified $\sigma,B$ of
VMA1 after fixing the original host data, its frame and comparison.
The arithmetic end now supplies \emph{one} absolute $g_*$ together
with $\mathrm{Marks}$, $\mathrm{FinEnd}$ and $\mathrm{AutEnd3Alt}$
at $(\sigma,g_*,B,K,\iota)$. Fix $g_1=p(g_*)$ before $R,u$,
the coordinate map or any representation. VMA1 descends the
marking. For each universal $R,u$ its weight-zero module is finite
by VMA2. For each subsequent coordinate point $r$ of Type I, take
the single integral representative $r_A$ just constructed in VMA5.
The proof of this realization uses only finiteness, independently
of any automorphy assertion. VMA1 applied to $r_A$ gives exactly
the $\mathrm{OrdPt}$ needed by $\mathrm{AutEnd3Alt}$. On every open
$U\subset H$, its matrices over $E$ are conjugate by $T(t)$ to
those of $r\circ p$, so VMA3 gives irreducibility there.

Choose an embedding of $K$ in $\Omega$ over $\mathbb Q$. Its
image is $\mathbb Q(i)$ and the image of its absolute Galois group
in $\Delta$ is $H$, since the action on $i$ defines $\eta_4$.
Let $\rho$ and $\rho_A$ be the resulting ordinary
$\mathrm{GL}_3(E)$ representations of $G_K$ obtained from
$r\circ p$ and $\mathrm{push}(r_A)\circ p$. They have literally
the specified matrices on restriction; their signs there are
positive. The conjugation $T(t)$ intertwines them, and still
intertwines their twists by $\chi_{K,\ell}^{-1}$.
Apply $\mathrm{AutEnd3Alt}$ to the absolute integral point and
$\rho_A$ to get $\mathrm{IsAutCsdA}$ for that twisted representation.
Transport this \emph{full} assertion to the twist of $\rho$.
For the first conjunct, the same regular algebraic
conjugate-self-dual cusp works: conjugation carries invariant
complements and preserves every local Frobenius polynomial of the
attachment. For the second conjunct, continuity is invariant under a
fixed conjugation; carrying the local flags retains their characters,
hence both $\mathrm{Unl1}$ and $\mathrm{PotUnl}$. For every cusp
satisfying the stated attachment, use the same real $t_0$ and the
same places in its $\mathrm{WinAt}$ assertion. This gives the
complete conditional $\mathrm{WitA}$ of the original Alt conclusion.
The chosen $K$-embedding, its scalar tower, $\rho$, and its literal
matrix identification therefore witness $\mathrm{AutOnAlt}(K,\iota,-1)$
for $r\circ p$. Apply VMA5 with this finite module and automorphy
at the unchanged mark to obtain $\mathrm{ThmS6}_s$.

What remains to obtain the whole selected assertion by this route is
the simultaneous supply of $\mathrm{FinEnd}$ and $\mathrm{AutEnd3Alt}$
at that mark for every original host. In $\mathrm{FinEnd}$ the
finite-index subgroup precedes all residual DVR points. In
$\mathrm{AutEnd3Alt}$ the conclusion is the inverse-cyclotomic
$\mathrm{IsAutCsdA}$ for every integral ordinary point satisfying the
open-subgroup hypothesis, at the fixed $\iota$. AFE16 and AFE21
supply an arithmetic frame and the implication from its rank-three
$\mathrm{FinPt}$ to $\mathrm{FinEnd}$ at their given cohomological
input. The pinned arithmetic route to the joint end also calls upon
the linked automorphic end, the rank-three local comparison at even
level, and conjugate-self-dual descent. VMA1--6 establish the
comparison from the joint end to the literal selected assertion;
the joint arithmetic end remains an input.
% END VMA6
